\subsubsection{关联压缩的等价性}
把服务区作为顶点，同一运输架次所涉及的区域构成一个关联集合；同一逻辑中继任务的实际保障关系再产生相应关联。通过任务节点构成二部图，或连接各关联集合内的成员，都可以得到相同的连通分量。连通闭包处理的是“必须同组”的传递关系，不是地理距离聚类。

必要性来自同组关系的传递性：区域$i$与$j$、$j$与$k$分别必须同组时，三者就须整体归组。因此合法分区不能切开连通分量。反之，将每个分量整体放入一个组后，所有运输和中继关联均留在组内，非空组要求满足即形成合法分区。原服务区合法分组与分量的合法分区于是等价。

当前15个服务区得到3个不可拆块，允许对两组的3种无标签分区和三组的唯一分区完整比较。压缩没有移动任务或删减需求，只消除了会破坏冻结执行关系的划分。原共享日程中同一飞机先后服务不同块，并不要求这些块同组：独立执行允许给各组重新配置同型实体，必须完整保留的是运输和逻辑中继任务本身。
